\documentclass[12pt,a4paper]{article}

\usepackage[T2A]{fontenc}
\usepackage[utf8]{inputenc}
\usepackage[russian]{babel}
\usepackage{amsmath,amssymb,amsthm,mathtools}
\usepackage{enumitem}
\usepackage{array}
\usepackage{booktabs}
\usepackage{hyperref}
\usepackage{listings}
\usepackage{makecell}
\usepackage{xcolor}

\hypersetup{
    colorlinks=true,
    linkcolor=blue,
    urlcolor=blue
}

\newtheorem{definition}{Определение}
\newtheorem{theorem}{Теорема}
\newtheorem{lemma}{Лемма}
\newtheorem{example}{Пример}

\DeclareMathOperator{\Enc}{Enc}
\DeclareMathOperator{\Dec}{Dec}
\DeclareMathOperator{\Gen}{Gen}
\DeclareMathOperator{\Sign}{Sign}
\DeclareMathOperator{\Ver}{Ver}
\DeclareMathOperator{\Mac}{Mac}
\DeclareMathOperator{\KeyGen}{KeyGen}
\DeclareMathOperator{\Hash}{Hash}
\DeclareMathOperator{\Adv}{Adv}

\title{Билет №17 \\ \small Вычислительная геометрия на плоскости. Уравнения точек, прямых, окружностей. Выпуклые оболочки, алгоритмы построения. Алгоритмы триангуляции. Задачи регионального поиска и локализации. Алгоритмы планирования движения. (ИТМО)}
\author{Сериков Владислав Александрович}
\date{5 мая 2026}

\begin{document}

\maketitle
\tableofcontents
\newpage

\section{Вычислительная геометрия на плоскости}

\subsection{Базовые объекты и обозначения}

\subsubsection{Точки и векторы}

Точка на плоскости задаётся парой координат
\[
p = (x_p, y_p) \in \mathbb{R}^2.
\]

Если даны две точки \(a=(x_a,y_a)\) и \(b=(x_b,y_b)\), то вектор из \(a\) в \(b\) равен
\[
\overrightarrow{ab} = b-a = (x_b-x_a,\; y_b-y_a).
\]

Длина вектора \(v=(x,y)\):
\[
\|v\| = \sqrt{x^2+y^2}.
\]

Расстояние между точками:
\[
\rho(a,b)=\|b-a\|=\sqrt{(x_b-x_a)^2+(y_b-y_a)^2}.
\]

\subsubsection{Скалярное произведение}

Для векторов \(u=(u_x,u_y)\), \(v=(v_x,v_y)\) скалярное произведение:
\[
\langle u,v\rangle = u_xv_x+u_yv_y.
\]

Также
\[
\langle u,v\rangle = \|u\|\|v\|\cos\varphi,
\]
где \(\varphi\) — угол между векторами.

Отсюда:
\[
u \perp v \Longleftrightarrow \langle u,v\rangle=0.
\]

\subsubsection{Псевдоскалярное произведение}

Для векторов \(u=(u_x,u_y)\), \(v=(v_x,v_y)\) определим
\[
[u,v] = u_xv_y-u_yv_x.
\]

Это ориентированная площадь параллелограмма, построенного на \(u\) и \(v\).

Если \(a,b,c\) — точки, то
\[
[\overrightarrow{ab},\overrightarrow{ac}]
=
(x_b-x_a)(y_c-y_a)-(y_b-y_a)(x_c-x_a).
\]

Знак этого выражения определяет ориентацию тройки точек:

\[
\begin{cases}
[\overrightarrow{ab},\overrightarrow{ac}]>0, & \text{поворот } a\to b\to c \text{ левый},\\
[\overrightarrow{ab},\overrightarrow{ac}]<0, & \text{поворот } a\to b\to c \text{ правый},\\
[\overrightarrow{ab},\overrightarrow{ac}]=0, & \text{точки коллинеарны}.
\end{cases}
\]

Площадь треугольника \(abc\):
\[
S_{\triangle abc}
=
\frac12
\left|
[\overrightarrow{ab},\overrightarrow{ac}]
\right|.
\]

\subsection{Прямые на плоскости}

\subsubsection{Общее уравнение прямой}

Прямая на плоскости задаётся уравнением
\[
Ax+By+C=0,
\]
где \(A,B,C\in\mathbb{R}\), причём \((A,B)\neq(0,0)\).

Вектор
\[
n=(A,B)
\]
является нормальным вектором прямой.

Направляющий вектор прямой можно взять равным
\[
d=(-B,A).
\]

\subsubsection{Прямая через две точки}

Пусть \(p_1=(x_1,y_1)\), \(p_2=(x_2,y_2)\), \(p_1\neq p_2\). Тогда прямая через эти точки задаётся уравнением
\[
\begin{vmatrix}
x & y & 1\\
x_1 & y_1 & 1\\
x_2 & y_2 & 1
\end{vmatrix}
=0.
\]

В раскрытом виде:
\[
(y_1-y_2)x+(x_2-x_1)y+(x_1y_2-x_2y_1)=0.
\]

\subsubsection{Параметрическое уравнение прямой}

Если прямая проходит через точку \(p_0=(x_0,y_0)\) и имеет направляющий вектор \(d=(d_x,d_y)\), то
\[
p(t)=p_0+td,
\qquad t\in\mathbb{R}.
\]

В координатах:
\[
\begin{cases}
x=x_0+td_x,\\
y=y_0+td_y.
\end{cases}
\]

\subsubsection{Луч и отрезок}

Луч:
\[
p(t)=p_0+td,\qquad t\ge 0.
\]

Отрезок с концами \(a\) и \(b\):
\[
p(t)=a+t(b-a),\qquad 0\le t\le 1.
\]

\subsubsection{Расстояние от точки до прямой}

Пусть прямая задана уравнением
\[
Ax+By+C=0.
\]
Тогда расстояние от точки \(p=(x_0,y_0)\) до прямой:
\[
d(p,l)=
\frac{|Ax_0+By_0+C|}{\sqrt{A^2+B^2}}.
\]

\subsubsection{Расстояние от точки до отрезка}

Пусть отрезок имеет концы \(a,b\), а точка \(p\). Обозначим
\[
v=b-a.
\]
Рассмотрим проекцию точки \(p\) на прямую \(ab\):
\[
t=\frac{\langle p-a,v\rangle}{\langle v,v\rangle}.
\]

Тогда ближайшая к \(p\) точка отрезка:
\[
q=
\begin{cases}
a, & t<0,\\
a+tv, & 0\le t\le 1,\\
b, & t>1.
\end{cases}
\]

Искомое расстояние:
\[
d(p,[ab])=\rho(p,q).
\]

\subsection{Пересечения прямых и отрезков}

\subsubsection{Пересечение двух прямых}

Пусть прямые заданы уравнениями
\[
A_1x+B_1y+C_1=0,
\]
\[
A_2x+B_2y+C_2=0.
\]

Рассмотрим определитель
\[
\Delta =
\begin{vmatrix}
A_1 & B_1\\
A_2 & B_2
\end{vmatrix}
=
A_1B_2-A_2B_1.
\]

Если \(\Delta\neq 0\), то прямые пересекаются в единственной точке:
\[
x=
\frac{
\begin{vmatrix}
-C_1 & B_1\\
-C_2 & B_2
\end{vmatrix}
}{\Delta}
=
\frac{B_1C_2-B_2C_1}{\Delta},
\]
\[
y=
\frac{
\begin{vmatrix}
A_1 & -C_1\\
A_2 & -C_2
\end{vmatrix}
}{\Delta}
=
\frac{C_1A_2-C_2A_1}{\Delta}.
\]

Если \(\Delta=0\), то прямые параллельны или совпадают.

\subsubsection{Критерий пересечения двух отрезков}

Пусть даны два отрезка \([ab]\) и \([cd]\).

Определим функцию ориентации:
\[
\operatorname{orient}(a,b,c)
=
[\overrightarrow{ab},\overrightarrow{ac}].
\]

\textbf{Общий случай.}
Если никакие три точки не коллинеарны, то отрезки пересекаются тогда и только тогда, когда
\[
\operatorname{orient}(a,b,c)\cdot \operatorname{orient}(a,b,d)<0
\]
и
\[
\operatorname{orient}(c,d,a)\cdot \operatorname{orient}(c,d,b)<0.
\]

\textbf{Полный случай с коллинеарностью.}
Точка \(p\) лежит на отрезке \([ab]\), если
\[
\operatorname{orient}(a,b,p)=0
\]
и
\[
\min(x_a,x_b)\le x_p\le \max(x_a,x_b),
\]
\[
\min(y_a,y_b)\le y_p\le \max(y_a,y_b).
\]

Отрезки пересекаются, если выполнено хотя бы одно из условий:
\[
\operatorname{orient}(a,b,c)\cdot \operatorname{orient}(a,b,d)<0
\quad\text{и}\quad
\operatorname{orient}(c,d,a)\cdot \operatorname{orient}(c,d,b)<0,
\]
или одна из точек \(a,b,c,d\) лежит на противоположном отрезке.

\subsubsection{Минимальный пример}

Пусть
\[
a=(0,0),\quad b=(4,4),
\]
\[
c=(0,4),\quad d=(4,0).
\]

Тогда
\[
\operatorname{orient}(a,b,c)
=
(4-0)(4-0)-(4-0)(0-0)=16>0,
\]
\[
\operatorname{orient}(a,b,d)
=
(4-0)(0-0)-(4-0)(4-0)=-16<0.
\]

Аналогично
\[
\operatorname{orient}(c,d,a)<0,
\qquad
\operatorname{orient}(c,d,b)>0.
\]

Следовательно, отрезки пересекаются.

\subsection{Окружности}

\subsubsection{Каноническое уравнение окружности}

Окружность с центром \(O=(a,b)\) и радиусом \(r>0\) задаётся уравнением
\[
(x-a)^2+(y-b)^2=r^2.
\]

\subsubsection{Общее уравнение окружности}

После раскрытия скобок:
\[
x^2+y^2+Dx+Ey+F=0.
\]

Это уравнение задаёт окружность, если
\[
\frac{D^2+E^2}{4}-F>0.
\]

Центр:
\[
O=
\left(
-\frac D2,\,
-\frac E2
\right).
\]

Радиус:
\[
r=
\sqrt{
\frac{D^2+E^2}{4}-F
}.
\]

\subsubsection{Окружность через три точки}

Пусть даны три неколлинеарные точки
\[
p_i=(x_i,y_i),\qquad i=1,2,3.
\]

Окружность через них можно найти из системы
\[
x_i^2+y_i^2+Dx_i+Ey_i+F=0,
\qquad i=1,2,3.
\]

Искомые \(D,E,F\) определяются решением линейной системы:
\[
\begin{cases}
D x_1+E y_1+F=-(x_1^2+y_1^2),\\
D x_2+E y_2+F=-(x_2^2+y_2^2),\\
D x_3+E y_3+F=-(x_3^2+y_3^2).
\end{cases}
\]

\subsubsection{Положение точки относительно окружности}

Пусть окружность имеет центр \(O\) и радиус \(r\). Для точки \(p\):

\[
\begin{cases}
\rho(p,O)<r, & p \text{ внутри окружности},\\
\rho(p,O)=r, & p \text{ на окружности},\\
\rho(p,O)>r, & p \text{ вне окружности}.
\end{cases}
\]

На практике часто сравнивают квадраты:
\[
(x_p-a)^2+(y_p-b)^2
\quad\text{и}\quad
r^2.
\]

\subsection{Ориентация и инцидентность}

\subsubsection{Теорема об ориентации трёх точек}

\textbf{Теорема.}
Пусть \(a,b,c\in\mathbb{R}^2\). Тогда знак определителя
\[
\Delta(a,b,c)
=
\begin{vmatrix}
x_a & y_a & 1\\
x_b & y_b & 1\\
x_c & y_c & 1
\end{vmatrix}
\]
определяет взаимное расположение точек:
\[
\Delta(a,b,c)>0
\]
тогда и только тогда, когда при движении от \(a\) к \(b\), а затем к \(c\), совершается левый поворот;
\[
\Delta(a,b,c)<0
\]
тогда и только тогда, когда совершается правый поворот;
\[
\Delta(a,b,c)=0
\]
тогда и только тогда, когда точки \(a,b,c\) коллинеарны.

\textbf{Доказательство.}
Вычтем первую строку определителя из второй и третьей:
\[
\Delta(a,b,c)
=
\begin{vmatrix}
x_a & y_a & 1\\
x_b-x_a & y_b-y_a & 0\\
x_c-x_a & y_c-y_a & 0
\end{vmatrix}.
\]

Разложим по третьему столбцу:
\[
\Delta(a,b,c)
=
\begin{vmatrix}
x_b-x_a & y_b-y_a\\
x_c-x_a & y_c-y_a
\end{vmatrix}.
\]

Значит,
\[
\Delta(a,b,c)
=
(x_b-x_a)(y_c-y_a)
-
(y_b-y_a)(x_c-x_a).
\]

Это псевдоскалярное произведение
\[
[\overrightarrow{ab},\overrightarrow{ac}].
\]

Его модуль равен площади параллелограмма, построенного на векторах
\(\overrightarrow{ab}\) и \(\overrightarrow{ac}\), а знак показывает ориентацию перехода от первого вектора ко второму. Положительный знак соответствует повороту против часовой стрелки, то есть левому повороту; отрицательный — повороту по часовой стрелке. Если определитель равен нулю, площадь параллелограмма равна нулю, следовательно, векторы линейно зависимы, а точки коллинеарны. Теорема доказана.
\[
\Box
\]

\section{Выпуклые множества и выпуклые оболочки}

\subsection{Выпуклые множества}

Множество \(S\subseteq\mathbb{R}^2\) называется выпуклым, если для любых двух точек \(a,b\in S\) весь отрезок \([ab]\) содержится в \(S\):
\[
\forall a,b\in S,\quad
\forall t\in[0,1]:
\quad
(1-t)a+tb\in S.
\]

\subsection{Выпуклая оболочка}

Выпуклой оболочкой множества точек \(P\) называется наименьшее выпуклое множество, содержащее \(P\).

Обозначение:
\[
\operatorname{conv}(P).
\]

Эквивалентно:
\[
\operatorname{conv}(P)
=
\left\{
\sum_{i=1}^{n}\lambda_i p_i
\;\middle|\;
p_i\in P,\;
\lambda_i\ge 0,\;
\sum_{i=1}^{n}\lambda_i=1
\right\}.
\]

Если \(P\) конечно и не все точки коллинеарны, то \(\operatorname{conv}(P)\) является выпуклым многоугольником.

\subsection{Опорные прямые}

Прямая \(l\) называется опорной прямой множества \(S\), если:
\[
S
\]
лежит целиком в одной из двух замкнутых полуплоскостей, определяемых \(l\), и
\[
l\cap S\neq\varnothing.
\]

Вершины выпуклой оболочки — это точки множества, через которые можно провести опорные прямые.

\subsection{Теорема о выпуклой оболочке конечного множества}

\textbf{Теорема.}
Пусть \(P\subset\mathbb{R}^2\) — конечное множество точек. Тогда \(\operatorname{conv}(P)\) является пересечением всех выпуклых множеств, содержащих \(P\), и одновременно множеством всех выпуклых комбинаций точек из \(P\).

\textbf{Доказательство.}

Обозначим через \(\mathcal{C}\) семейство всех выпуклых множеств, содержащих \(P\):
\[
\mathcal{C}=\{C\subseteq\mathbb{R}^2\mid C \text{ выпукло и } P\subseteq C\}.
\]

Рассмотрим множество
\[
K=\bigcap_{C\in\mathcal{C}} C.
\]

Поскольку \(P\subseteq C\) для любого \(C\in\mathcal{C}\), то
\[
P\subseteq K.
\]

Пересечение произвольного семейства выпуклых множеств выпукло. Действительно, если \(a,b\in K\), то \(a,b\in C\) для любого \(C\in\mathcal{C}\). Так как каждый \(C\) выпукл, то
\[
(1-t)a+tb\in C
\]
для всех \(t\in[0,1]\) и всех \(C\in\mathcal{C}\). Следовательно,
\[
(1-t)a+tb\in K.
\]
Значит, \(K\) выпукло.

Поскольку \(K\) выпукло и содержит \(P\), оно является одним из выпуклых множеств, содержащих \(P\), причём оно содержится в любом таком множестве. Поэтому \(K\) — наименьшее выпуклое множество, содержащее \(P\), то есть
\[
K=\operatorname{conv}(P).
\]

Теперь рассмотрим множество всех выпуклых комбинаций точек из \(P\):
\[
H=
\left\{
\sum_{i=1}^{n}\lambda_i p_i
\;\middle|\;
p_i\in P,\;
\lambda_i\ge 0,\;
\sum_{i=1}^{n}\lambda_i=1
\right\}.
\]

Покажем, что \(H\) выпукло. Пусть
\[
x=\sum_i\lambda_i p_i,
\qquad
y=\sum_j\mu_j q_j,
\]
где \(p_i,q_j\in P\),
\[
\lambda_i,\mu_j\ge 0,
\qquad
\sum_i\lambda_i=1,
\qquad
\sum_j\mu_j=1.
\]

Для \(t\in[0,1]\):
\[
(1-t)x+ty
=
\sum_i (1-t)\lambda_i p_i
+
\sum_j t\mu_j q_j.
\]

Все коэффициенты неотрицательны, и их сумма равна
\[
(1-t)\sum_i\lambda_i+t\sum_j\mu_j
=
(1-t)+t=1.
\]

Значит, \((1-t)x+ty\) является выпуклой комбинацией точек из \(P\), то есть принадлежит \(H\). Следовательно, \(H\) выпукло.

Очевидно, \(P\subseteq H\), поскольку каждая точка \(p\in P\) представляется как выпуклая комбинация с коэффициентом \(1\).

Так как \(H\) выпукло и содержит \(P\), то
\[
\operatorname{conv}(P)\subseteq H.
\]

С другой стороны, любое выпуклое множество, содержащее \(P\), содержит все отрезки между точками из \(P\), затем все выпуклые комбинации этих точек. Поэтому
\[
H\subseteq \operatorname{conv}(P).
\]

Следовательно,
\[
H=\operatorname{conv}(P).
\]

Теорема доказана.
\[
\Box
\]

\section{Алгоритмы построения выпуклой оболочки}

\subsection{Алгоритм Джарвиса}

Алгоритм Джарвиса также называют \emph{заворачиванием подарка}. Он строит выпуклую оболочку, последовательно находя следующую вершину оболочки.

\subsubsection{Идея}

Начинаем с самой левой точки. Затем выбираем такую следующую точку, что все остальные точки лежат справа от направленного ребра текущая точка \(\to\) следующая точка. Это эквивалентно выбору самой ``левой'' или самой ``внешней'' точки относительно текущей.

\subsubsection{Алгоритм}

Пусть дано множество точек
\[
P=\{p_1,\dots,p_n\}.
\]

\begin{enumerate}
    \item Найти стартовую точку \(p_0\) с минимальной координатой \(x\). При равенстве выбрать точку с минимальной координатой \(y\).
    \item Положить \(h_0=p_0\).
    \item Для текущей вершины \(h_i\) выбрать произвольную точку \(q\neq h_i\).
    \item Для каждой точки \(r\in P\) проверить ориентацию тройки \((h_i,q,r)\).
    \item Если
    \[
    \operatorname{orient}(h_i,q,r)<0,
    \]
    то заменить \(q\) на \(r\).
    \item После просмотра всех точек \(q\) является следующей вершиной оболочки:
    \[
    h_{i+1}=q.
    \]
    \item Повторять, пока не вернёмся в \(p_0\).
\end{enumerate}

\subsubsection{Сложность}

Если \(h\) — число вершин выпуклой оболочки, то для каждой вершины выполняется просмотр всех \(n\) точек. Поэтому сложность:
\[
O(nh).
\]

В худшем случае \(h=n\), поэтому:
\[
O(n^2).
\]

\subsubsection{Минимальный пример}

Пусть
\[
P=\{(0,0),(1,1),(2,0),(1,2),(1,0)\}.
\]

Схематично:

\[
\begin{array}{ccc}
 & (1,2) & \\
(0,0) & (1,1) & (2,0)\\
 & (1,0) &
\end{array}
\]

Выпуклая оболочка:
\[
(0,0)\to(2,0)\to(1,2)\to(0,0).
\]

Точка \((1,1)\) и \((1,0)\) лежат внутри или на границе.

\subsection{Алгоритм Грэхема}

\subsubsection{Идея}

Алгоритм Грэхема сортирует точки по полярному углу относительно специальной стартовой точки, затем проходит по отсортированному списку, поддерживая стек вершин текущей оболочки.

\subsubsection{Алгоритм}

\begin{enumerate}
    \item Найти точку \(p_0\) с минимальной координатой \(y\). При равенстве выбрать точку с минимальной координатой \(x\).
    \item Отсортировать остальные точки по полярному углу относительно \(p_0\).
    \item При равных углах оставить наиболее удалённую от \(p_0\) точку.
    \item Поместить первые две точки в стек.
    \item Последовательно обрабатывать оставшиеся точки.
    \item Пока последние две точки стека \(a,b\) и текущая точка \(c\) образуют не левый поворот, то есть
    \[
    \operatorname{orient}(a,b,c)\le 0,
    \]
    удалить \(b\) из стека.
    \item Добавить \(c\) в стек.
    \item После обработки всех точек стек содержит вершины выпуклой оболочки в порядке обхода.
\end{enumerate}

\subsubsection{Псевдокод}

\[
\begin{array}{l}
\textsc{GrahamScan}(P):\\
\quad p_0 \leftarrow \text{точка с минимальным } y\\
\quad \text{отсортировать } P\setminus\{p_0\} \text{ по полярному углу вокруг } p_0\\
\quad S \leftarrow \text{пустой стек}\\
\quad \text{push}(S,p_0)\\
\quad \text{push}(S,p_1)\\
\quad \text{push}(S,p_2)\\
\quad \text{для } i=3,\dots,n:\\
\quad\quad \text{пока } |S|\ge 2 \text{ и } \operatorname{orient}(\text{next-to-top}(S),\text{top}(S),p_i)\le 0:\\
\quad\quad\quad \text{pop}(S)\\
\quad\quad \text{push}(S,p_i)\\
\quad \text{вернуть } S
\end{array}
\]

\subsubsection{Сложность}

Сортировка:
\[
O(n\log n).
\]

Каждая точка добавляется в стек один раз и удаляется не более одного раза. Поэтому проход по стеку:
\[
O(n).
\]

Итого:
\[
O(n\log n).
\]

\subsubsection{Пример}

Пусть точки:
\[
(0,0),(1,1),(2,0),(2,2),(0,2),(1,0).
\]

Стартовая точка:
\[
p_0=(0,0).
\]

После сортировки по углу:
\[
(1,0),(2,0),(1,1),(2,2),(0,2).
\]

Точки с одинаковым углом относительно \((0,0)\):
\[
(1,0),(2,0)
\]
имеют угол \(0\). Оставляем более дальнюю:
\[
(2,0).
\]

Итоговая оболочка:
\[
(0,0)\to(2,0)\to(2,2)\to(0,2)\to(0,0).
\]

\subsection{Алгоритм Эндрю}

Алгоритм Эндрю часто называют монотонной цепью. Он строит нижнюю и верхнюю части выпуклой оболочки.

\subsubsection{Алгоритм}

\begin{enumerate}
    \item Отсортировать точки лексикографически:
    \[
    (x_1,y_1)<(x_2,y_2)
    \Longleftrightarrow
    x_1<x_2
    \text{ или }
    (x_1=x_2 \text{ и } y_1<y_2).
    \]
    \item Построить нижнюю оболочку:
    \begin{enumerate}
        \item Идти по точкам слева направо.
        \item Поддерживать стек.
        \item Пока последние две точки стека и текущая точка не образуют левый поворот, удалять верхнюю точку стека.
        \item Добавить текущую точку.
    \end{enumerate}
    \item Построить верхнюю оболочку аналогично, проходя точки справа налево.
    \item Объединить нижнюю и верхнюю оболочки, удалив дублирующиеся крайние точки.
\end{enumerate}

\subsubsection{Псевдокод}

\[
\begin{array}{l}
\textsc{AndrewHull}(P):\\
\quad \text{отсортировать } P \text{ лексикографически}\\
\quad L \leftarrow []\\
\quad \text{для } p\in P:\\
\quad\quad \text{пока } |L|\ge 2 \text{ и } \operatorname{orient}(L[-2],L[-1],p)\le 0:\\
\quad\quad\quad \text{удалить } L[-1]\\
\quad\quad \text{добавить } p \text{ в } L\\
\\
\quad U \leftarrow []\\
\quad \text{для } p\in \text{reverse}(P):\\
\quad\quad \text{пока } |U|\ge 2 \text{ и } \operatorname{orient}(U[-2],U[-1],p)\le 0:\\
\quad\quad\quad \text{удалить } U[-1]\\
\quad\quad \text{добавить } p \text{ в } U\\
\\
\quad \text{вернуть } L[:-1] + U[:-1]
\end{array}
\]

\subsubsection{Сложность}

Сортировка:
\[
O(n\log n).
\]

Построение двух цепей:
\[
O(n).
\]

Итого:
\[
O(n\log n).
\]

\subsubsection{Минимальный пример}

Пусть
\[
P=\{(0,0),(1,1),(2,0),(2,2),(0,2)\}.
\]

После сортировки:
\[
(0,0),(0,2),(1,1),(2,0),(2,2).
\]

Нижняя оболочка:
\[
(0,0)\to(2,0)\to(2,2).
\]

Верхняя оболочка:
\[
(2,2)\to(0,2)\to(0,0).
\]

Итог:
\[
(0,0)\to(2,0)\to(2,2)\to(0,2).
\]

\subsection{Нижняя оценка для построения выпуклой оболочки}

\textbf{Теорема.}
В модели вычислений, основанной на сравнениях, построение выпуклой оболочки \(n\) точек на плоскости в худшем случае требует
\[
\Omega(n\log n)
\]
операций.

\textbf{Доказательство.}
Сведём задачу сортировки \(n\) вещественных чисел к задаче построения выпуклой оболочки.

Пусть даны числа
\[
x_1,x_2,\dots,x_n.
\]

Построим точки
\[
p_i=(x_i,x_i^2).
\]

Все эти точки лежат на параболе
\[
y=x^2.
\]

Парабола является строго выпуклой кривой. Если числа \(x_i\) различны, то все точки \(p_i\) являются вершинами нижней части выпуклой оболочки.

Действительно, для любых трёх точек на параболе с
\[
x_i<x_j<x_k
\]
точка \(p_j\) лежит ниже хорды \(p_ip_k\). Поэтому при обходе нижней оболочки точки идут в порядке возрастания \(x_i\).

Следовательно, если бы существовал алгоритм построения выпуклой оболочки быстрее, чем
\[
\Omega(n\log n),
\]
то можно было бы отсортировать числа \(x_1,\dots,x_n\) быстрее, чем за
\[
\Omega(n\log n)
\]
сравнений.

Но в модели сравнений сортировка требует
\[
\Omega(n\log n)
\]
операций в худшем случае. Получаем противоречие.

Значит, построение выпуклой оболочки требует
\[
\Omega(n\log n).
\]

Теорема доказана.
\[
\Box
\]

\section{Триангуляция}

\subsection{Определение триангуляции}

Пусть \(P\) — простой многоугольник. Триангуляцией многоугольника называется разбиение его внутренности на непересекающиеся треугольники с вершинами в вершинах многоугольника.

Диагональю простого многоугольника называется отрезок, соединяющий две несоседние вершины многоугольника и полностью лежащий внутри многоугольника, возможно кроме концов.

\subsection{Теорема о существовании триангуляции}

\textbf{Теорема.}
Любой простой многоугольник с \(n\ge 3\) вершинами допускает триангуляцию.

\textbf{Доказательство.}
Докажем по индукции по числу вершин \(n\).

База: при \(n=3\) многоугольник уже является треугольником, поэтому триангуляция существует.

Пусть утверждение верно для всех простых многоугольников с числом вершин меньше \(n\). Рассмотрим простой многоугольник \(P\) с \(n>3\) вершинами.

Известно, что в любом простом многоугольнике существует хотя бы одна диагональ. Докажем это.

Возьмём вершину \(v\), имеющую минимальную \(x\)-координату. Пусть её соседние вершины — \(u\) и \(w\). Поскольку \(v\) — самая левая вершина, внутренний угол при \(v\) меньше \(\pi\) или больше \(\pi\), но треугольник \(\triangle uvw\) расположен с правой стороны от \(v\).

Если отрезок \(uw\) целиком лежит внутри многоугольника, то \(uw\) является диагональю.

Если внутри треугольника \(\triangle uvw\) лежат вершины многоугольника, выберем среди них вершину \(z\), наиболее удалённую от прямой \(uw\) в сторону \(v\). Тогда отрезок \(vz\) лежит внутри многоугольника и является диагональю. Если бы он пересекал границу многоугольника, то существовала бы вершина внутри треугольника, расположенная ещё дальше, что противоречит выбору \(z\).

Итак, существует диагональ. Она разбивает многоугольник \(P\) на два простых многоугольника \(P_1\) и \(P_2\), каждый из которых имеет меньше \(n\) вершин. По предположению индукции оба они триангулируемы.

Объединяя их триангуляции, получаем триангуляцию исходного многоугольника \(P\).

Теорема доказана.
\[
\Box
\]

\subsection{Число треугольников в триангуляции}

\textbf{Теорема.}
Любая триангуляция простого многоугольника с \(n\) вершинами содержит ровно
\[
n-2
\]
треугольника и
\[
n-3
\]
диагонали.

\textbf{Доказательство.}
Пусть в триангуляции:
\[
T
\]
треугольников,
\[
D
\]
диагоналей.

Каждый треугольник имеет \(3\) стороны. Всего по всем треугольникам получаем \(3T\) сторон.

Стороны исходного многоугольника учитываются один раз. Их \(n\). Каждая диагональ является общей стороной двух треугольников и поэтому учитывается дважды. Следовательно,
\[
3T = n+2D.
\]

С другой стороны, каждая новая диагональ увеличивает число областей внутри многоугольника на \(1\). Чтобы получить \(T\) треугольников, нужно провести
\[
D=T-1
\]
диагоналей.

Подставим:
\[
3T=n+2(T-1).
\]

Тогда:
\[
3T=n+2T-2,
\]
\[
T=n-2.
\]

Следовательно:
\[
D=T-1=n-3.
\]

Теорема доказана.
\[
\Box
\]

\subsection{Алгоритм отсечения ушей}

\subsubsection{Ухо многоугольника}

Пусть \(P\) — простой многоугольник. Три последовательные вершины
\[
v_{i-1},v_i,v_{i+1}
\]
образуют ухо, если:

\begin{enumerate}
    \item вершина \(v_i\) выпуклая:
    \[
    \operatorname{orient}(v_{i-1},v_i,v_{i+1})>0
    \]
    при обходе многоугольника против часовой стрелки;
    \item диагональ \([v_{i-1}v_{i+1}]\) лежит внутри многоугольника;
    \item внутри треугольника
    \[
    \triangle v_{i-1}v_iv_{i+1}
    \]
    нет других вершин многоугольника.
\end{enumerate}

\subsubsection{Теорема о двух ушах}

\textbf{Теорема.}
У любого простого многоугольника с \(n\ge 4\) вершинами существуют по крайней мере два непересекающихся уха.

\textbf{Доказательство.}
Поскольку любой простой многоугольник триангулируем, рассмотрим произвольную его триангуляцию.

Триангуляция многоугольника с \(n\) вершинами содержит
\[
n-2
\]
треугольника и
\[
n-3
\]
диагонали.

Построим двойственный граф триангуляции: каждой треугольной области сопоставим вершину графа, а две вершины двойственного графа соединим ребром, если соответствующие треугольники имеют общую диагональ.

Поскольку диагонали не пересекаются и разбиение происходит внутри простого многоугольника, двойственный граф является деревом.

В любом дереве с более чем одной вершиной существуют как минимум два листа.

Лист двойственного графа соответствует треугольнику, который имеет общую диагональ только с одним соседним треугольником. Значит, две другие его стороны являются сторонами исходного многоугольника. Такой треугольник имеет вид
\[
\triangle v_{i-1}v_iv_{i+1}.
\]

Следовательно, вершина \(v_i\) образует ухо. Так как в дереве не менее двух листьев, то в многоугольнике есть по крайней мере два уха.

Теорема доказана.
\[
\Box
\]

\subsubsection{Алгоритм}

\begin{enumerate}
    \item Задать многоугольник списком вершин в порядке обхода против часовой стрелки:
    \[
    v_1,v_2,\dots,v_n.
    \]
    \item Пока число вершин больше \(3\):
    \begin{enumerate}
        \item Найти вершину \(v_i\), образующую ухо.
        \item Добавить треугольник
        \[
        \triangle v_{i-1}v_iv_{i+1}
        \]
        в ответ.
        \item Удалить вершину \(v_i\) из многоугольника.
    \end{enumerate}
    \item Оставшиеся три вершины образуют последний треугольник.
\end{enumerate}

\subsubsection{Проверка уха}

Для вершины \(v_i\):

\begin{enumerate}
    \item Проверить выпуклость:
    \[
    \operatorname{orient}(v_{i-1},v_i,v_{i+1})>0.
    \]
    \item Проверить, что никакая другая вершина многоугольника не лежит внутри треугольника
    \[
    \triangle v_{i-1}v_iv_{i+1}.
    \]
\end{enumerate}

Точка \(p\) лежит внутри треугольника \(abc\), если ориентации имеют одинаковые знаки:
\[
\operatorname{orient}(a,b,p)\ge 0,
\]
\[
\operatorname{orient}(b,c,p)\ge 0,
\]
\[
\operatorname{orient}(c,a,p)\ge 0
\]
для обхода против часовой стрелки.

\subsubsection{Сложность}

Наивная реализация:
\[
O(n^3),
\]
поскольку для каждого удаления можно искать ухо за \(O(n^2)\).

Более аккуратная реализация с поддержанием списка ушей:
\[
O(n^2).
\]

\subsubsection{Минимальный пример}

Пусть многоугольник задан вершинами:
\[
(0,0),(3,0),(3,2),(1,1),(0,2).
\]

Схема:
\[
\begin{array}{ccccc}
(0,2) & & & (3,2)\\
 & (1,1) & & \\
(0,0) & & & (3,0)
\end{array}
\]

Одно из ушей:
\[
(0,0),(3,0),(3,2).
\]

После его отсечения остаётся четырёхугольник:
\[
(0,0),(3,2),(1,1),(0,2).
\]

Затем отсекается следующее ухо, пока не останется треугольник.

\subsection{Триангуляция монотонного многоугольника}

\subsubsection{Монотонный многоугольник}

Многоугольник называется монотонным относительно оси \(y\), если любая горизонтальная прямая пересекает его границу не более чем в двух точках.

Эквивалентно: его граница состоит из двух \(y\)-монотонных цепей, соединяющих верхнюю и нижнюю вершины.

\subsubsection{Алгоритм триангуляции монотонного многоугольника}

\begin{enumerate}
    \item Отсортировать вершины по убыванию \(y\)-координаты.
    \item Определить, к какой цепи относится каждая вершина: левой или правой.
    \item Поместить первые две вершины в стек.
    \item Последовательно обрабатывать вершины \(v_i\).
    \item Если \(v_i\) лежит на другой цепи, чем вершина на вершине стека:
    \begin{enumerate}
        \item соединить \(v_i\) диагоналями со всеми вершинами стека, кроме последней;
        \item очистить стек;
        \item положить в стек предыдущую вершину и \(v_i\).
    \end{enumerate}
    \item Если \(v_i\) лежит на той же цепи:
    \begin{enumerate}
        \item пока диагональ от \(v_i\) к предпоследней вершине стека лежит внутри многоугольника, проводить её и удалять верхнюю вершину;
        \item добавить \(v_i\) в стек.
    \end{enumerate}
    \item После обработки всех вершин провести оставшиеся диагонали к вершинам стека.
\end{enumerate}

\subsubsection{Сложность}

Сортировка:
\[
O(n\log n).
\]

Если вершины уже отсортированы вдоль границы, монотонный многоугольник можно триангулировать за
\[
O(n).
\]

\subsection{Триангуляция Делоне}

\subsubsection{Определение}

Пусть \(P\) — множество точек на плоскости. Триангуляция множества \(P\) называется триангуляцией Делоне, если для любого треугольника её описанная окружность не содержит внутри себя других точек из \(P\).

Формально, для каждого треугольника
\[
\triangle abc
\]
и каждой точки
\[
p\in P\setminus\{a,b,c\}
\]
выполнено:
\[
p
\]
лежит вне или на границе окружности, описанной около
\[
\triangle abc.
\]

\subsubsection{Локальный критерий Делоне}

Пусть два треугольника
\[
\triangle abc
\quad\text{и}\quad
\triangle abd
\]
имеют общее ребро \(ab\) и образуют выпуклый четырёхугольник \(acbd\).

Ребро \(ab\) является локально делонеевским, если точка \(d\) не лежит внутри окружности, описанной около \(\triangle abc\), и точка \(c\) не лежит внутри окружности, описанной около \(\triangle abd\).

Если ребро не удовлетворяет этому условию, его можно заменить на диагональ \(cd\). Эта операция называется \emph{флипом ребра}.

\subsubsection{Тест принадлежности точки окружности}

Для ориентированного против часовой стрелки треугольника \(abc\) точка \(p\) лежит внутри описанной окружности тогда и только тогда, когда
\[
\begin{vmatrix}
x_a & y_a & x_a^2+y_a^2 & 1\\
x_b & y_b & x_b^2+y_b^2 & 1\\
x_c & y_c & x_c^2+y_c^2 & 1\\
x_p & y_p & x_p^2+y_p^2 & 1
\end{vmatrix}
>0.
\]

Если определитель равен нулю, точка лежит на окружности.

\subsubsection{Инкрементальный алгоритм построения триангуляции Делоне}

\begin{enumerate}
    \item Построить большой вспомогательный треугольник, содержащий все точки \(P\).
    \item Начать с триангуляции, состоящей только из этого треугольника.
    \item Последовательно добавлять точки из \(P\).
    \item Для очередной точки \(p\):
    \begin{enumerate}
        \item найти треугольник, содержащий \(p\);
        \item разбить его на три треугольника, соединив \(p\) с вершинами найденного треугольника;
        \item проверить новые рёбра на условие Делоне;
        \item если ребро не локально делонеевское, выполнить флип;
        \item рекурсивно проверить соседние рёбра.
    \end{enumerate}
    \item Удалить все треугольники, содержащие вершины вспомогательного большого треугольника.
\end{enumerate}

\subsubsection{Сложность}

При случайном порядке вставки и эффективной локализации точки ожидаемая сложность:
\[
O(n\log n).
\]

В худшем случае при неудачной реализации:
\[
O(n^2).
\]

\subsubsection{Минимальный пример}

Пусть
\[
P=\{(0,0),(1,0),(0,1),(1,1)\}.
\]

Возможны две триангуляции квадрата:
\[
(0,0)-(1,1)
\]
или
\[
(1,0)-(0,1).
\]

Если все четыре точки лежат на одной окружности, обе триангуляции являются делонеевскими, поскольку ни одна точка не лежит строго внутри описанной окружности другого треугольника.

\section{Региональный поиск}

\subsection{Постановка задачи}

Региональный поиск — это задача нахождения всех объектов из заданного множества, попадающих в некоторую область запроса.

Чаще всего рассматриваются точки:
\[
P=\{p_1,\dots,p_n\}\subset\mathbb{R}^2.
\]

Запрос задаёт область \(R\), и нужно найти
\[
P\cap R.
\]

Примеры областей запроса:

\begin{itemize}
    \item прямоугольник со сторонами, параллельными осям;
    \item круг;
    \item полуплоскость;
    \item произвольный многоугольник.
\end{itemize}

\subsection{Одномерное дерево поиска}

Для одномерного множества чисел
\[
X=\{x_1,\dots,x_n\}
\]
запрос:
\[
[a,b].
\]

Если хранить числа в сбалансированном бинарном дереве поиска, то можно найти все числа из диапазона \([a,b]\) за
\[
O(\log n+k),
\]
где \(k\) — число найденных элементов.

\subsection{\texorpdfstring{\(k\)-d}{kd}-дерево}

\subsubsection{Идея}

\(k\)-d-дерево — это бинарное дерево, рекурсивно разбивающее пространство гиперплоскостями, перпендикулярными координатным осям.

В двумерном случае:

\begin{itemize}
    \item на первом уровне делим по \(x\);
    \item на втором уровне делим по \(y\);
    \item затем снова по \(x\), и так далее.
\end{itemize}

\subsubsection{Построение двумерного \(k\)-d-дерева}

\begin{enumerate}
    \item Если множество точек пусто, вернуть пустое дерево.
    \item Выбрать координату разбиения:
    \[
    x, y, x, y,\dots
    \]
    в зависимости от глубины.
    \item Найти медиану точек по выбранной координате.
    \item Поместить медианную точку в текущий узел.
    \item Точки с меньшей координатой отправить в левое поддерево.
    \item Точки с большей координатой отправить в правое поддерево.
    \item Рекурсивно построить поддеревья.
\end{enumerate}

\subsubsection{Запрос прямоугольником}

Пусть запрос:
\[
R=[x_{\min},x_{\max}]\times[y_{\min},y_{\max}].
\]

Алгоритм:

\begin{enumerate}
    \item В узле дерева проверить, лежит ли точка узла внутри \(R\).
    \item Если лежит, добавить её в ответ.
    \item В зависимости от разделяющей координаты определить, какие поддеревья могут пересекать \(R\).
    \item Рекурсивно обойти только эти поддеревья.
\end{enumerate}

\subsubsection{Сложность}

Для сбалансированного двумерного \(k\)-d-дерева запрос прямоугольником выполняется за
\[
O(\sqrt n+k)
\]
в худшем случае.

Память:
\[
O(n).
\]

\subsubsection{Минимальный пример}

Пусть точки:
\[
(1,2),(2,3),(3,1),(4,4),(5,2).
\]

На первом уровне делим по \(x\). Медиана:
\[
(3,1).
\]

Левое поддерево:
\[
(1,2),(2,3).
\]

Правое поддерево:
\[
(4,4),(5,2).
\]

Запрос:
\[
R=[1,4]\times[2,4].
\]

Ответ:
\[
(1,2),(2,3),(4,4).
\]

\subsection{Range tree}

\subsubsection{Идея}

Двумерное range tree — это дерево поиска по одной координате, в каждой вершине которого хранится вспомогательная структура по второй координате.

\subsubsection{Построение}

\begin{enumerate}
    \item Построить сбалансированное дерево поиска по координате \(x\).
    \item Для каждой вершины \(v\) хранить список всех точек её поддерева, отсортированный по координате \(y\).
\end{enumerate}

\subsubsection{Запрос прямоугольником}

Для запроса
\[
R=[x_1,x_2]\times[y_1,y_2]
\]
нужно:

\begin{enumerate}
    \item В основном дереве найти канонические поддеревья, покрывающие все точки с
    \[
    x\in[x_1,x_2].
    \]
    Их число:
    \[
    O(\log n).
    \]
    \item В каждом найденном поддереве выполнить одномерный поиск по списку \(y\)-координат:
    \[
    y\in[y_1,y_2].
    \]
\end{enumerate}

\subsubsection{Сложность}

Без fractional cascading:
\[
O(\log^2 n+k).
\]

С fractional cascading:
\[
O(\log n+k).
\]

Память:
\[
O(n\log n).
\]

\subsection{Quad-tree}

\subsubsection{Идея}

Quad-tree рекурсивно разбивает квадратную область на четыре квадранта.

\subsubsection{Построение}

\begin{enumerate}
    \item Взять квадрат, содержащий все точки.
    \item Если в текущей ячейке мало точек или достигнута максимальная глубина, сделать её листом.
    \item Иначе разделить квадрат на четыре равных квадрата.
    \item Разнести точки по дочерним квадратам.
    \item Рекурсивно повторить.
\end{enumerate}

\subsubsection{Запрос}

Для области \(R\):

\begin{enumerate}
    \item Начать с корневой ячейки.
    \item Если ячейка не пересекает \(R\), пропустить её.
    \item Если ячейка целиком лежит внутри \(R\), добавить все точки ячейки.
    \item Иначе рекурсивно обработать дочерние ячейки.
\end{enumerate}

Quad-tree хорошо работает для пространственно равномерных данных, но в худшем случае может иметь большую глубину.

\section{Локализация точки}

\subsection{Постановка задачи}

Дано разбиение плоскости на области, например планарное подразделение:
\[
\mathcal{S}.
\]

Для заданной точки \(q\) нужно определить область, содержащую \(q\).

Такая задача называется локализацией точки.

Примеры:

\begin{itemize}
    \item определить треугольник триангуляции, содержащий точку;
    \item определить грань планарного графа;
    \item определить регион карты.
\end{itemize}

\subsection{Наивный алгоритм}

Если разбиение состоит из \(m\) областей, можно проверить принадлежность точки каждой области.

Для треугольной сетки:

\begin{enumerate}
    \item Для каждого треугольника \(\triangle abc\).
    \item Проверить, лежит ли \(q\) внутри него.
    \item Первый найденный треугольник вернуть.
\end{enumerate}

Сложность:
\[
O(m).
\]

\subsection{Локализация в триангуляции методом обхода}

Если дана триангуляция и для каждого треугольника известны соседи по рёбрам, можно использовать алгоритм ``walking''.

\subsubsection{Алгоритм}

\begin{enumerate}
    \item Начать с некоторого треугольника \(T\).
    \item Если точка \(q\) лежит в \(T\), вернуть \(T\).
    \item Иначе найти ребро \(e\), через которое луч от некоторой точки внутри \(T\) к \(q\) выходит из \(T\).
    \item Перейти в соседний треугольник через ребро \(e\).
    \item Повторять до нахождения треугольника.
\end{enumerate}

В худшем случае:
\[
O(n).
\]

На практике для хорошо построенных триангуляций часто работает быстрее.

\subsection{Трапецоидальная карта}

\subsubsection{Идея}

Для множества непересекающихся отрезков строится вертикальная декомпозиция: из концов отрезков проводятся вертикальные лучи вверх и вниз до ближайшего отрезка или границы.

Полученные области называются трапециями, хотя некоторые могут быть треугольниками или неограниченными областями.

\subsubsection{Структура поиска}

Для локализации точки строится ориентированный ациклический граф поиска. В его внутренних вершинах находятся вопросы двух типов:

\begin{itemize}
    \item сравнение \(x_q\) с \(x\)-координатой некоторой точки;
    \item проверка, лежит ли \(q\) выше или ниже некоторого отрезка.
\end{itemize}

Листья соответствуют трапециям.

\subsubsection{Инкрементальное построение}

\begin{enumerate}
    \item Начать с большой ограничивающей рамки.
    \item Вставлять отрезки в случайном порядке.
    \item Для нового отрезка найти последовательность трапеций, которые он пересекает.
    \item Разбить эти трапеции новым отрезком и вертикальными лучами из его концов.
    \item Обновить граф поиска.
\end{enumerate}

\subsubsection{Сложность}

Для случайного порядка вставки ожидаемые оценки:

\[
O(n\log n)
\]
на построение,

\[
O(n)
\]
памяти,

\[
O(\log n)
\]
на запрос локализации.

\section{Планирование движения}

\subsection{Постановка задачи}

В задачах планирования движения имеется робот, препятствия и начальное и конечное положения.

В простейшем случае:

\begin{itemize}
    \item робот — точка на плоскости;
    \item препятствия — многоугольники;
    \item требуется найти непрерывный путь от стартовой точки \(s\) до целевой точки \(t\), не пересекающий препятствий.
\end{itemize}

Свободное пространство:
\[
F=\mathbb{R}^2\setminus O,
\]
где \(O\) — объединение препятствий.

Задача: найти непрерывную кривую
\[
\gamma:[0,1]\to F
\]
такую, что
\[
\gamma(0)=s,
\qquad
\gamma(1)=t.
\]

\subsection{Конфигурационное пространство}

Если робот не является точкой, вместо физического пространства рассматривают пространство конфигураций.

Конфигурация — это набор параметров, однозначно задающий положение робота.

Например:

\begin{itemize}
    \item точечный робот на плоскости:
    \[
    q=(x,y);
    \]
    \item диск радиуса \(r\):
    \[
    q=(x,y);
    \]
    \item жёсткое тело на плоскости:
    \[
    q=(x,y,\theta).
    \]
\end{itemize}

Препятствия в конфигурационном пространстве:
\[
C_{\text{obs}}.
\]

Свободное конфигурационное пространство:
\[
C_{\text{free}}=C\setminus C_{\text{obs}}.
\]

Планирование движения сводится к поиску пути в
\[
C_{\text{free}}.
\]

\subsection{Робот-диск}

Если робот — диск радиуса \(r\), то задачу можно свести к задаче точечного робота, расширив препятствия на радиус \(r\).

Для препятствия \(O\) расширенное препятствие:
\[
O\oplus B_r
=
\{o+b\mid o\in O,\; b\in B_r\},
\]
где
\[
B_r=\{x\in\mathbb{R}^2\mid \|x\|\le r\}.
\]

Операция \(\oplus\) называется суммой Минковского.

Тогда центр диска должен двигаться в пространстве:
\[
\mathbb{R}^2\setminus \bigcup_i (O_i\oplus B_r).
\]

\subsection{Граф видимости}

\subsubsection{Идея}

Для точечного робота среди многоугольных препятствий кратчайший путь, если он существует, является ломаной, вершины которой принадлежат множеству:
\[
\{s,t\}\cup\{\text{вершины препятствий}\}.
\]

Граф видимости строится на этих вершинах. Две вершины соединяются ребром, если отрезок между ними не пересекает внутренности препятствий.

\subsubsection{Теорема о кратчайшем пути}

\textbf{Теорема.}
Пусть точечный робот движется на плоскости среди многоугольных препятствий. Если между точками \(s\) и \(t\) существует кратчайший путь в свободном пространстве, то этот путь является ломаной, промежуточные вершины которой являются вершинами препятствий.

\textbf{Доказательство.}
Рассмотрим кратчайший путь \(\gamma\) от \(s\) до \(t\).

Если некоторый участок пути лежит во внутренности свободного пространства и не касается препятствий, то кратчайшее соединение его концов — отрезок. Поэтому такой участок можно заменить прямолинейным отрезком, не увеличив длину пути.

Следовательно, изменение направления пути может происходить только в точках, где путь касается препятствий.

Пусть путь меняет направление в точке \(p\), лежащей на ребре препятствия, но не являющейся его вершиной. Тогда в достаточно малой окрестности точки \(p\) препятствие выглядит как полуплоскость, ограниченная прямой. Две соседние части пути подходят к \(p\) и выходят из \(p\) в одной и той же свободной полуплоскости. В этом случае можно заменить две малые части пути на более короткий отрезок, лежащий в той же свободной полуплоскости. Это противоречит минимальности пути.

Значит, повороты кратчайшего пути могут происходить только в вершинах препятствий.

Таким образом, кратчайший путь является ломаной с вершинами среди
\[
s,t
\]
и вершин препятствий.

Теорема доказана.
\[
\Box
\]

\subsubsection{Алгоритм построения графа видимости}

Пусть множество вершин:
\[
V=\{s,t\}\cup\{\text{вершины препятствий}\}.
\]

\begin{enumerate}
    \item Создать пустой граф
    \[
    G=(V,E).
    \]
    \item Для каждой пары вершин \(u,v\in V\):
    \begin{enumerate}
        \item построить отрезок \([uv]\);
        \item проверить, пересекает ли \([uv]\) внутренность какого-либо препятствия;
        \item если не пересекает, добавить ребро
        \[
        (u,v)
        \]
        с весом
        \[
        \rho(u,v).
        \]
    \end{enumerate}
    \item Запустить алгоритм Дейкстры из \(s\) в \(t\).
\end{enumerate}

\subsubsection{Сложность}

Если всего \(n\) вершин препятствий, то пар вершин:
\[
O(n^2).
\]

Наивная проверка видимости одной пары может занимать
\[
O(n).
\]

Итого построение:
\[
O(n^3).
\]

Существуют более эффективные алгоритмы построения графа видимости за
\[
O(n^2\log n).
\]

\subsubsection{Минимальный пример}

Пусть есть одно прямоугольное препятствие:
\[
[1,3]\times[1,2],
\]
старт
\[
s=(0,0),
\]
цель
\[
t=(4,0).
\]

Прямой отрезок \([st]\) не пересекает препятствие, поэтому кратчайший путь:
\[
(0,0)\to(4,0).
\]

Если цель:
\[
t=(4,1.5),
\]
то прямой путь пересекает препятствие. Тогда путь идёт через одну из видимых вершин прямоугольника, например:
\[
(0,0)\to(1,1)\to(3,1)\to(4,1.5),
\]
если соответствующие отрезки лежат в свободном пространстве.

\subsection{Вероятностные дорожные карты}

Метод вероятностных дорожных карт, или PRM, применяется в высокоразмерных конфигурационных пространствах.

\subsubsection{Алгоритм PRM}

\begin{enumerate}
    \item Сгенерировать случайные конфигурации
    \[
    q_1,\dots,q_m\in C_{\text{free}}.
    \]
    \item Добавить стартовую и целевую конфигурации:
    \[
    s,t.
    \]
    \item Для каждой конфигурации найти несколько ближайших соседей.
    \item Попытаться соединить соседние конфигурации локальным планировщиком, обычно прямолинейным движением в конфигурационном пространстве.
    \item Если локальный путь не пересекает препятствий, добавить ребро в граф.
    \item Найти путь в полученном графе от \(s\) до \(t\).
\end{enumerate}

\subsubsection{Свойства}

PRM не гарантирует нахождения пути за конечное число случайных проб, но обладает вероятностной полнотой: если путь существует и имеет положительный зазор от препятствий, то вероятность найти путь стремится к \(1\) при увеличении числа сэмплов.

\subsection{RRT}

RRT, Rapidly-exploring Random Tree, строит дерево, быстро исследующее свободное пространство.

\subsubsection{Алгоритм RRT}

\begin{enumerate}
    \item Инициализировать дерево:
    \[
    T=\{s\}.
    \]
    \item Повторять:
    \begin{enumerate}
        \item выбрать случайную конфигурацию
        \[
        q_{\text{rand}};
        \]
        \item найти ближайшую вершину дерева
        \[
        q_{\text{near}};
        \]
        \item сделать шаг от \(q_{\text{near}}\) в сторону \(q_{\text{rand}}\):
        \[
        q_{\text{new}};
        \]
        \item если движение от \(q_{\text{near}}\) к \(q_{\text{new}}\) не сталкивается с препятствиями, добавить
        \[
        q_{\text{new}}
        \]
        и ребро
        \[
        (q_{\text{near}},q_{\text{new}})
        \]
        в дерево;
        \item если \(q_{\text{new}}\) достаточно близко к цели, попытаться соединить его с \(t\).
    \end{enumerate}
\end{enumerate}

\subsubsection{Особенности}

RRT хорошо работает в пространствах большой размерности и быстро исследует ранее не покрытые области. Однако обычный RRT не оптимален по длине пути.

\subsection{RRT*}

RRT* модифицирует RRT, добавляя переподключение соседних вершин для улучшения стоимости пути.

\subsubsection{Основные шаги}

\begin{enumerate}
    \item Сгенерировать случайную конфигурацию.
    \item Найти ближайшую вершину.
    \item Добавить новую вершину.
    \item Рассмотреть множество соседних вершин в некотором радиусе.
    \item Выбрать родителя, дающего минимальную стоимость пути от старта.
    \item Переподключить соседей через новую вершину, если это уменьшает их стоимость.
\end{enumerate}

RRT* является асимптотически оптимальным: при увеличении числа итераций стоимость найденного пути стремится к оптимальной.

\section{Численная устойчивость в вычислительной геометрии}

\subsection{Проблема ошибок округления}

Многие геометрические алгоритмы опираются на знаки определителей:
\[
\operatorname{orient}(a,b,c)
\]
или тест принадлежности окружности.

При вычислениях с плавающей точкой возможны ошибки, особенно когда точки почти коллинеарны или почти лежат на одной окружности.

\subsection{Подходы}

\begin{itemize}
    \item использовать целочисленную арифметику, если координаты целые;
    \item использовать длинные целые типы для определителей;
    \item применять точную рациональную арифметику;
    \item применять адаптивные предикаты, которые используют точные вычисления только в трудных случаях.
\end{itemize}

Например, если координаты целые и по модулю не превосходят \(M\), то ориентация:
\[
(x_b-x_a)(y_c-y_a)-(y_b-y_a)(x_c-x_a)
\]
может иметь порядок
\[
M^2.
\]

Поэтому при больших координатах тип \texttt{int} может переполниться, и следует использовать \texttt{long long} или произвольную точность.

\section{Итоговая таблица алгоритмов}

\[
\begin{array}{|l|l|l|}
\hline
\text{Задача} & \text{Алгоритм} & \text{Сложность}\\
\hline
\text{Выпуклая оболочка} & \text{Джарвис} & O(nh)\\
\hline
\text{Выпуклая оболочка} & \text{Грэхем} & O(n\log n)\\
\hline
\text{Выпуклая оболочка} & \text{Эндрю} & O(n\log n)\\
\hline
\makecell[l]{\text{Триангуляция простого} \\ \text{многоугольника}} & \text{Отсечение ушей} & O(n^2)\text{--}O(n^3)\\
\hline
\makecell[l]{\text{Триангуляция монотонного} \\ \text{многоугольника}} & \text{Стековый алгоритм} & \makecell[l]{O(n) \text{ при} \\ \text{отсортированных вершинах}}\\
\hline
\text{Триангуляция Делоне} & \text{Инкрементальный алгоритм} & O(n\log n)\text{ ожидаемо}\\
\hline
\text{Региональный поиск} & k\text{-d-дерево} & O(\sqrt n+k)\\
\hline
\text{Региональный поиск} & \text{Range tree} & O(\log^2 n+k)\\
\hline
\text{Локализация точки} & \text{Трапецоидальная карта} & O(\log n)\text{ ожидаемо}\\
\hline
\text{Планирование движения} & \text{Граф видимости} & O(n^3)\text{ наивно}\\
\hline
\text{Планирование движения} & \text{PRM, RRT, RRT*} & \text{зависит от числа сэмплов}\\
\hline
\end{array}
\]

\end{document}
